\begin{derivation}[从\eqnrefs{eq:dirac-mass-chiral-mixing-dp11}到\eqnrefs{eq:sm-yukawa-invariant-structures-dp43}]
以带电轻子为例，$L_L$ 与 $e_R$ 的表示分别为

\begin{equation*}
 L_L:(\mathbf1,\mathbf2,-\tfrac12),
 \qquad
 e_R:(\mathbf1,\mathbf1,-1).
\end{equation*}
共轭场 $\bar L_L$ 属于 $SU(2)$ 共轭二重态，超荷符号反转。因此裸双线性满足
\begin{equation*}
 \bar L_L e_R:
 \qquad
 \bar{\mathbf2}\otimes\mathbf1=\bar{\mathbf2},
 \qquad
 Y=+\frac12-1=-\frac12.
\end{equation*}
它既不是 $SU(2)_L$ 单态，也不是 $U(1)_Y$ 中性量，因而不能出现在未破缺电弱理论的局域拉氏量中。夸克左右手场具有同样的二重态--单态错配，所以裸夸克 Dirac 质量也被排除。

插入 Higgs 二重态后，带电轻子组合的弱同位旋部分为
\begin{equation*}
 \bar{\mathbf2}\otimes\mathbf2
 =\mathbf1\oplus\mathbf3,
\end{equation*}
其中包含单态；取指标收缩 $\bar L_{Li}\Phi_i$ 即选出这一单态。其总超荷为
\begin{equation*}
 -Y(L_L)+Y(\Phi)+Y(e_R)
 =+\frac12+\frac12-1=0,
\end{equation*}
故 $\bar L_L\Phi e_R$ 在全部电弱规范因子下不变。

对下型夸克，弱同位旋指标由 $\bar Q_{Li}\Phi_i$ 收缩成 $SU(2)_L$ 单态。其总超荷为
\begin{align*}
 Y(\bar Q_L\Phi d_R)
 &=-Y(Q_L)+Y(\Phi)+Y(d_R),\\
 &=-\frac16+\frac12-\frac13,\\
 &=0.
\end{align*}
颜色指标则由 $\bar{\mathbf3}\otimes\mathbf3$ 中的单态收缩，因此 $\bar Q_L\Phi d_R$ 在三个规范因子下都为单态。上型夸克若使用 $\Phi$，超荷不能抵消；改用
\begin{equation*}
 \widetilde\Phi:(\mathbf1,\mathbf2,-\tfrac12)
\end{equation*}
后，
\begin{equation*}
 -Y(Q_L)+Y(\widetilde\Phi)+Y(u_R)
 =-\frac16-\frac12+\frac23=0,
\end{equation*}
于是 $\bar Q_L\widetilde\Phi u_R$ 也是规范单态。这三种组合恰好给出\eqnrefs{eq:sm-yukawa-invariant-structures-dp43}。
\end{derivation}

